\documentclass[11pt,a4paper]{article}

% ---------- Packages ----------
\usepackage[margin=1in]{geometry}
\usepackage{amsmath, amssymb, amsthm, mathtools}
\usepackage{microtype}
\usepackage{algorithm}
\usepackage{algpseudocode}
\usepackage[colorlinks=true,
            linkcolor=blue!50!black,
            urlcolor=blue!50!black,
            citecolor=blue!50!black]{hyperref}

% ---------- Theorem environments ----------
\theoremstyle{plain}
\newtheorem{theorem}{Theorem}[section]
\newtheorem{proposition}[theorem]{Proposition}
\newtheorem{lemma}[theorem]{Lemma}
\newtheorem{corollary}[theorem]{Corollary}

\theoremstyle{definition}
\newtheorem{definition}[theorem]{Definition}

\theoremstyle{remark}
\newtheorem{remark}[theorem]{Remark}
\newtheorem{example}[theorem]{Example}

% ---------- Math shortcuts ----------
\newcommand{\R}{\mathbb{R}}
\newcommand{\N}{\mathbb{N}}
\newcommand{\E}{\mathbb{E}}
\newcommand{\Var}{\operatorname{Var}}
\newcommand{\Cov}{\operatorname{Cov}}
\newcommand{\Corr}{\operatorname{Corr}}
\newcommand{\F}{\mathcal{F}}
\newcommand{\Prob}{\mathbb{P}}
\newcommand{\Q}{\mathbb{Q}}
\newcommand{\dd}{\mathop{}\!\mathrm{d}}
\newcommand{\eqdef}{\overset{\mathrm{def}}{=}}
\newcommand{\indic}{\mathbf{1}}

% ---------- Document metadata ----------
\title{Greeks by Monte Carlo:\\
       Bumping, Pathwise, and Likelihood Ratio}
\author{Quant Finance Portfolio --- Phase 2, Block 4}
\date{\today}

\begin{document}
\maketitle

\begin{abstract}
We compute the European-call Greeks (Delta, Vega, Gamma) by Monte
Carlo using three increasingly sophisticated techniques: bumping
(finite differences applied to the pricer), pathwise sensitivities
(differentiation of the payoff under the integral), and likelihood
ratio (differentiation of the kernel of the underlying density). All
three are unbiased for the appropriate Greeks under appropriate
regularity conditions; they trade off variance, bias, and
applicability differently. We derive each from first principles,
implement them under the exact GBM sampler of Block~1.1, and
validate against closed-form Black-Scholes Greeks at the canonical
ATM point. The block also surfaces a delicate but important point:
the pathwise method does \emph{not} apply to Gamma (the kink of
$\max(\cdot, 0)$ produces a Dirac contribution after the second
derivative), and the likelihood-ratio method is known to be
high-variance for Gamma; in practice Gamma is computed by central
finite differences.
\end{abstract}

\tableofcontents

\section{The setting}
\label{sec:setting}

The price of a European call under risk-neutral GBM is
\begin{equation}
C(S_0, \sigma) \;=\; e^{-rT}\,\E\bigl[\max(S_T - K, 0)\bigr],
\qquad
S_T \;=\; S_0 \exp\bigl((r - \tfrac{1}{2}\sigma^2)T + \sigma\sqrt T\,Z\bigr),
\qquad Z \sim \mathcal{N}(0, 1).
\end{equation}
The Greeks of interest are
\begin{equation}
\Delta \;=\; \frac{\partial C}{\partial S_0},
\qquad
\nu \;=\; \frac{\partial C}{\partial \sigma},
\qquad
\Gamma \;=\; \frac{\partial^2 C}{\partial S_0^2}.
\end{equation}
Black-Scholes provides closed forms (Block~1):
\begin{equation}
\Delta^{\mathrm{BS}} = \Phi(d_1),
\qquad
\nu^{\mathrm{BS}} = S_0 \phi(d_1) \sqrt T,
\qquad
\Gamma^{\mathrm{BS}} = \frac{\phi(d_1)}{S_0 \sigma \sqrt T},
\end{equation}
with $\phi$ the standard normal pdf and $d_1 = (\log(S_0/K) + (r +
\sigma^2/2)T)/(\sigma\sqrt T)$. We use these as ground truth.

The MC pricer of Block~1.1 produces an estimate
$\hat C_n = \frac{1}{n}\sum_{i=1}^n e^{-rT}\max(S_T^{(i)} - K, 0)$
where $S_T^{(i)} = S_0 \exp((r-\frac{1}{2}\sigma^2)T + \sigma\sqrt T
Z_i)$ for $Z_i$ iid $\mathcal N(0,1)$. Greeks are MC estimators of
$\partial C / \partial \theta$ for $\theta \in \{S_0, \sigma\}$ or
of the second derivative.

\section{Bumping (finite differences)}
\label{sec:bumping}

The conceptually simplest Greek estimator: differentiate the
\emph{pricer}, not the payoff.

\subsection{Construction}

Let $\hat C_n(\theta)$ be the MC estimate at parameter $\theta$.
Choose a small bump $h$. The forward-difference estimator of
$\partial C/\partial\theta$ is
\begin{equation}
\widehat{\partial_\theta C}_n^{\mathrm{FD}}
\;=\; \frac{\hat C_n(\theta + h) - \hat C_n(\theta)}{h}.
\end{equation}
The central-difference estimator is
\begin{equation}
\widehat{\partial_\theta C}_n^{\mathrm{CD}}
\;=\; \frac{\hat C_n(\theta + h) - \hat C_n(\theta - h)}{2h}.
\end{equation}
For Gamma we apply central differences twice:
\begin{equation}
\widehat{\Gamma}_n
\;=\; \frac{\hat C_n(S_0 + h) - 2\hat C_n(S_0) + \hat C_n(S_0 - h)}{h^2}.
\end{equation}

\subsection{Common random numbers (CRN)}

Naively, the two evaluations $\hat C_n(\theta + h)$ and $\hat
C_n(\theta - h)$ would use independent Monte Carlo samples. The
variance of the difference is then $\Var(\hat C_n(\theta + h)) +
\Var(\hat C_n(\theta - h)) \approx 2 \sigma^2/n$, and dividing by
$2h$ amplifies this further. The result: a bump-and-revalue Greek
has variance scaling as $\sigma^2/(n h^2)$, which blows up as
$h \to 0$.

The standard fix is \emph{common random numbers}: use the same set
of $\{Z_i\}$ for both evaluations. Pathwise on each $i$, the only
thing that changes between the two evaluations is the parameter,
not the noise. The variance of the difference is then much smaller
because the noise mostly cancels. For our setting,
\begin{equation}
\Var\Bigl(\hat C_n(S_0 + h, \{Z_i\}) - \hat C_n(S_0 - h, \{Z_i\})\Bigr)
\;\approx\; h^2 \cdot \Var(\partial_{S_0} f_i)
\end{equation}
to leading order in $h$, where $f_i = e^{-rT}\max(S_T^{(i)} - K, 0)$.
The factor of $h^2$ cancels with the $1/h^2$ in $\Gamma$, leaving a
finite variance estimator.

\subsection{The bias-variance trade-off in $h$}

With CRN, the bumping estimator has:
\begin{itemize}
\item Bias: $O(h)$ for forward differences, $O(h^2)$ for central
differences.
\item Variance: roughly constant in $h$ (after CRN cancellation).
\end{itemize}
The optimal $h$ minimises bias$^2$ + variance. For central
differences with bias $O(h^2)$ and variance $O(1)$, the optimum is
where both contribute equally, giving $h_\star$ on the order of
$\epsilon^{1/4}$ where $\epsilon$ is the floating-point precision
($\sim 10^{-16}$), so $h_\star \sim 10^{-4}$ for double precision.
We use $h = 10^{-2}$ for stability against MC noise; the bias is
then $\sim 10^{-4}$, comfortably below the half-width at our
budgets.

For Gamma, the second-order central difference has bias $O(h^2)$
but variance $O(1/h^2)$ even with CRN (the leading-order
cancellation that worked for $\Delta$ does not extend to the second
derivative). Gamma is intrinsically harder.

\section{Pathwise sensitivities}
\label{sec:pathwise}

The bumping method differentiates the pricer. The pathwise method
differentiates the \emph{payoff}.

\subsection{The interchange of differentiation and expectation}

Suppose $f(\theta, Z)$ is the payoff (a function of the parameter
$\theta$ and the underlying randomness $Z$), and we want to
estimate $\partial_\theta \E[f(\theta, Z)]$. If $f$ is sufficiently
regular, we can interchange differentiation and expectation:
\begin{equation}
\frac{\partial}{\partial\theta} \E[f(\theta, Z)]
\;=\; \E\!\left[\frac{\partial f(\theta, Z)}{\partial \theta}\right].
\end{equation}
The Monte Carlo estimator of the right-hand side is
\begin{equation}
\widehat{\partial_\theta C}_n^{\mathrm{PW}}
\;=\; \frac{1}{n} \sum_{i=1}^n
\frac{\partial f(\theta, Z_i)}{\partial \theta}.
\end{equation}
This is unbiased for the true derivative if the interchange is
justified, which by dominated convergence requires $f$ Lipschitz in
$\theta$ uniformly in $Z$ (or weaker conditions). The European call
payoff $\max(S_T - K, 0)$ is Lipschitz in $S_T$ (with constant 1),
which is in turn $C^\infty$ in $S_0$ and $\sigma$, so the
interchange is valid.

\subsection{Pathwise Delta and Vega for the European call}

Compute $\partial f/\partial S_0$ first. Using
\begin{equation}
S_T = S_0 \exp\bigl((r - \tfrac{1}{2}\sigma^2)T + \sigma\sqrt T \,Z\bigr)
\;=\; \frac{S_T}{S_0} \cdot S_0,
\end{equation}
so $\partial S_T / \partial S_0 = S_T / S_0$. By the chain rule
(and noting that $\max(x, 0)$ has derivative $\indic_{\{x > 0\}}$
almost everywhere),
\begin{equation}
\frac{\partial f}{\partial S_0}
\;=\; e^{-rT} \cdot \indic_{\{S_T > K\}} \cdot \frac{S_T}{S_0}.
\end{equation}

For Vega:
\begin{equation}
\frac{\partial S_T}{\partial \sigma}
\;=\; S_T \cdot (-\sigma T + \sqrt T \, Z),
\end{equation}
so
\begin{equation}
\frac{\partial f}{\partial \sigma}
\;=\; e^{-rT} \cdot \indic_{\{S_T > K\}} \cdot S_T \cdot (\sqrt T\, Z - \sigma T).
\end{equation}

The pathwise Delta and Vega estimators are then
\begin{equation}
\widehat{\Delta}_n^{\mathrm{PW}}
\;=\; \frac{1}{n} \sum_{i=1}^n e^{-rT} \indic_{\{S_T^{(i)} > K\}}
\frac{S_T^{(i)}}{S_0},
\end{equation}
\begin{equation}
\widehat{\nu}_n^{\mathrm{PW}}
\;=\; \frac{1}{n} \sum_{i=1}^n e^{-rT} \indic_{\{S_T^{(i)} > K\}}
S_T^{(i)} (\sqrt T\, Z_i - \sigma T).
\end{equation}

These are unbiased and have variance comparable to (often less
than) the variance of the price estimator itself. They are the
gold standard when applicable.

\subsection{Why pathwise fails for Gamma}

For Gamma we want $\partial^2 C / \partial S_0^2$, which would
require differentiating $\partial f/\partial S_0$ once more.
Computing:
\begin{equation}
\frac{\partial^2 f}{\partial S_0^2}
\;=\; \frac{\partial}{\partial S_0}\bigl[
e^{-rT} \indic_{\{S_T > K\}} \frac{S_T}{S_0}\bigr].
\end{equation}
The indicator $\indic_{\{S_T > K\}}$ is the Heaviside step
function of $S_T - K$. Its derivative in the distributional sense
is the Dirac delta $\delta(S_T - K)$. Once we differentiate, the
expectation
\begin{equation}
\E\!\left[\delta(S_T - K) \cdot (\text{stuff})\right]
\end{equation}
is finite (it equals (stuff at $S_T = K$) times the density of $S_T$
at $K$), but \emph{not estimable by Monte Carlo}: a single sample
has $S_T = K$ with probability zero, so the sum is identically zero
on every realisation.

This is the canonical failure of pathwise for second-order Greeks
of payoffs with kinks: the second derivative is a measure, not a
function, and Monte Carlo can only estimate functions.

\section{Likelihood ratio}
\label{sec:lr}

The likelihood-ratio (LR) method, developed by Broadie and
Glasserman, sidesteps the problem of differentiating the payoff by
differentiating the \emph{density of the underlying} instead.

\subsection{The construction}

Suppose $S_T$ has a density $p(s; \theta)$ depending on the parameter
$\theta$. Then
\begin{equation}
C(\theta) \;=\; e^{-rT} \int f(s)\, p(s; \theta)\, \dd s,
\end{equation}
and (assuming we can interchange differentiation and integration)
\begin{equation}
\frac{\partial C}{\partial\theta}
\;=\; e^{-rT} \int f(s) \cdot \frac{\partial p(s; \theta)/\partial\theta}{p(s; \theta)}
\cdot p(s; \theta)\, \dd s
\;=\; e^{-rT} \E\!\left[ f(S_T) \cdot \mathcal{S}(\theta, S_T) \right],
\end{equation}
where
\begin{equation}
\mathcal{S}(\theta, S_T)
\;\eqdef\; \frac{\partial}{\partial\theta} \log p(S_T; \theta)
\end{equation}
is the \emph{score function} associated with the density family.

The MC estimator is
\begin{equation}
\widehat{\partial_\theta C}_n^{\mathrm{LR}}
\;=\; \frac{1}{n} \sum_{i=1}^n e^{-rT}\, f(S_T^{(i)}) \cdot \mathcal{S}(\theta, S_T^{(i)}).
\end{equation}

\subsection{Score function for GBM}

Under GBM, $\log S_T = \log S_0 + (r - \frac{1}{2}\sigma^2)T +
\sigma\sqrt T \, Z$ is normal with mean $\mu = \log S_0 + (r -
\frac{1}{2}\sigma^2)T$ and standard deviation $\nu = \sigma\sqrt T$.
The density of $\log S_T$ is
\begin{equation}
p_{\log}(\ell; S_0, \sigma)
\;=\; \frac{1}{\nu \sqrt{2\pi}}
\exp\!\left(-\frac{(\ell - \mu)^2}{2\nu^2}\right).
\end{equation}
Differentiating $\log p_{\log}$ with respect to $S_0$:
\begin{equation}
\mathcal{S}_{S_0}(\log S_T) \;=\; \frac{Z}{S_0 \sigma \sqrt T}.
\end{equation}
With respect to $\sigma$:
\begin{equation}
\mathcal{S}_\sigma(\log S_T)
\;=\; \frac{Z^2 - 1}{\sigma} - Z \sqrt T.
\end{equation}
(Both derivations use the chain rule and the fact that $Z = (\log
S_T - \mu)/\nu$.)

The LR Delta and Vega estimators are
\begin{equation}
\widehat{\Delta}_n^{\mathrm{LR}}
\;=\; \frac{1}{n} \sum_{i=1}^n e^{-rT} \max(S_T^{(i)} - K, 0)
\cdot \frac{Z_i}{S_0 \sigma \sqrt T},
\end{equation}
\begin{equation}
\widehat{\nu}_n^{\mathrm{LR}}
\;=\; \frac{1}{n} \sum_{i=1}^n e^{-rT} \max(S_T^{(i)} - K, 0)
\cdot \left(\frac{Z_i^2 - 1}{\sigma} - Z_i \sqrt T\right).
\end{equation}

\subsection{Variance comparison: pathwise vs LR}

For the European call, both pathwise and LR are unbiased. The
variance comparison is:
\begin{itemize}
\item Pathwise: the estimand is $e^{-rT}\indic_{\{S_T>K\}}(S_T/S_0)$,
which is bounded for typical paths and has moderate variance.
\item LR: the estimand is $f(S_T) \cdot Z / (S_0 \sigma \sqrt T)$,
which carries the full payoff $f$ multiplied by a factor that can
have large absolute value when $|Z|$ is large. Tails of $Z$
contribute $f(S_T) \cdot Z$, which can be large.
\end{itemize}
Empirically, for the call, pathwise has variance roughly $2$--$3
\times$ smaller than LR. The advantage of LR is universality:
for digital options where $f$ is a Heaviside (not Lipschitz),
pathwise fails but LR works.

\section{Algorithmic specifications}
\label{sec:algorithms}

\begin{algorithm}[H]
\caption{Bump-and-revalue Delta with central differences and CRN}
\label{alg:bump}
\begin{algorithmic}[1]
\Require Spot $S_0$, parameters, sample size $n$, bump $h$, rng.
\State Generate $\{Z_i\}_{i=1}^n$ iid $\mathcal N(0,1)$.
\State $C^+ \gets \frac{1}{n} \sum_i e^{-rT} \max(S_0(1+h) e^{(r-\sigma^2/2)T + \sigma\sqrt T Z_i} - K, 0)$
\State $C^- \gets \frac{1}{n} \sum_i e^{-rT} \max(S_0(1-h) e^{(r-\sigma^2/2)T + \sigma\sqrt T Z_i} - K, 0)$
\State $\Delta \gets (C^+ - C^-)/(2 S_0 h)$
\State \Return $\Delta$
\end{algorithmic}
\end{algorithm}

The trick: we bump $S_0$ multiplicatively (by $1 \pm h$) so that
the formulae work uniformly without recomputing the diffusion
coefficient.

\begin{algorithm}[H]
\caption{Pathwise Delta}
\label{alg:pathwise}
\begin{algorithmic}[1]
\Require Spot $S_0$, parameters, sample size $n$, rng.
\State Generate $\{Z_i\}_{i=1}^n$ iid $\mathcal N(0, 1)$.
\For{$i = 1, \ldots, n$}
\State $S_T^{(i)} \gets S_0 \exp((r - \tfrac{1}{2}\sigma^2)T + \sigma\sqrt T Z_i)$
\State $\Delta_i \gets e^{-rT} \cdot \indic_{\{S_T^{(i)} > K\}} \cdot S_T^{(i)} / S_0$
\EndFor
\State \Return $\bar\Delta = n^{-1} \sum_i \Delta_i$
\end{algorithmic}
\end{algorithm}

\begin{algorithm}[H]
\caption{Likelihood-ratio Delta}
\label{alg:lr}
\begin{algorithmic}[1]
\Require Spot $S_0$, parameters, sample size $n$, rng.
\State Generate $\{Z_i\}_{i=1}^n$ iid $\mathcal N(0, 1)$.
\For{$i = 1, \ldots, n$}
\State $S_T^{(i)} \gets S_0 \exp((r - \tfrac{1}{2}\sigma^2)T + \sigma\sqrt T Z_i)$
\State $\Delta_i \gets e^{-rT} \max(S_T^{(i)} - K, 0) \cdot Z_i / (S_0 \sigma \sqrt T)$
\EndFor
\State \Return $\bar\Delta = n^{-1} \sum_i \Delta_i$
\end{algorithmic}
\end{algorithm}

\section{Empirical validation strategy}
\label{sec:validation}

The script \texttt{python/validate\_greeks\_mc.py} runs four
tests.

\paragraph{Test 1: BS coherence.} For each Greek and each method,
the estimator at $n = 10^5$ should agree with the closed-form BS
Greek within a few standard errors.

\paragraph{Test 2: variance comparison Delta.} Pathwise vs LR vs
bumping. Predicted: pathwise $<$ LR $<$ bumping. Report variance
ratios at fixed budget.

\paragraph{Test 3: variance comparison Vega.} Same as Test 2.

\paragraph{Test 4: input validation.} Mechanical: bad inputs
should raise.

The Block~4 implementation is a single new module
\texttt{quantlib.greeks} containing nine pricers (3 methods
$\times$ 3 Greeks, with the caveat that pathwise/LR Gamma is
``not implemented'' and falls back to bumping).

\section{Bibliographic notes}
\label{sec:biblio}

The pathwise method for option Greeks is due to Broadie and
Glasserman~\cite{BroadieGlasserman1996}. The likelihood-ratio
method is also from \cite{BroadieGlasserman1996}; the broader
context of efficient simulation of derivatives is in
\cite[Chapter~7]{Glasserman}. The malliavin calculus generalisation,
which provides a unified framework for both pathwise and LR, is in
Fourni\'e et al.~\cite{FournieEtAl1999}. We do not pursue the
Malliavin generalisation in this writeup but note that it is the
modern theoretical lens.

\begin{thebibliography}{99}

\bibitem{BroadieGlasserman1996}
M.\ Broadie and P.\ Glasserman,
\emph{Estimating security price derivatives using simulation},
Management Science, 42(2):269--285, 1996.

\bibitem{FournieEtAl1999}
E.\ Fourni\'e, J.\ M.\ Lasry, J.\ Lebuchoux, P.\ L.\ Lions, N.\ Touzi,
\emph{Applications of Malliavin calculus to Monte-Carlo methods in
finance},
Finance and Stochastics, 3(4):391--412, 1999.

\bibitem{Glasserman}
P.\ Glasserman,
\emph{Monte Carlo Methods in Financial Engineering},
Springer, 2004.

\end{thebibliography}

\end{document}
